%% Motivic status of the Virasoro residue tower.

\providecommand{\MZV}{\operatorname{MZV}}
\providecommand{\mot}{\mathfrak m}
\providecommand{\Vir}{\mathrm{Vir}}
\providecommand{\cA}{\mathcal A}
\providecommand{\per}{\operatorname{per}}
\providecommand{\ClaimStatusProvedHere}{}
\providecommand{\ClaimStatusConditional}{}
\providecommand{\ClaimStatusDefinitional}{}
\providecommand{\ClaimStatusOpen}{}

\chapter{Motivic Structure of the Virasoro Residue Tower}
\label{ch:virasoro-motivic-purity-all-r-platonic}

\section*{The comparison problem}

The Virasoro OPE is defined over \(\mathbb Q(c)\).  This arithmetic
fact governs the local Ward identities.  A higher scalar shadow is a
second construction: it uses an ordered residue complex, a contraction
onto a scalar line, and the homotopies that transfer the Maurer--Cartan
equation.  A motivic statement requires one further datum, namely a
motivic lift of that residue construction.  The chapter records the
exact implication supplied by these data and isolates the remaining
comparison theorem.

We use four named packages:
\[
 \mathsf H_{\mathrm{Vir}},\qquad
 \mathsf H_{\mathrm{res}},\qquad
 \mathsf H_{\mathrm{mot}},\qquad
 \mathsf H_{\mathrm{Tate}}.
\]
Here \(\mathsf H_{\mathrm{Vir}}\) is the local Virasoro vacuum package,
\(\mathsf H_{\mathrm{res}}\) is the ordered residue package of
Chapter~\ref{ch:shadow-tower-higher-coefficients},
\(\mathsf H_{\mathrm{mot}}\) is a motivic realization of its chains and
homotopies, and \(\mathsf H_{\mathrm{Tate}}\) asserts that the final
scalar projection factors through the Tate line \(\mathbb Q(c)(0)\).

\section{Motivic periods and residue data}
\label{sec:vmpar-setup}

\begin{definition}[Motivic multiple-zeta algebra]
\label{def:vmpar-mzv-ring}
\ClaimStatusDefinitional
Let \(\mathcal H^{\mot}\) denote Brown's graded algebra of motivic
multiple zeta values over \(\mathbb Z\).  Its weight-zero component is
\(\mathcal H^{\mot}_0=\mathbb Q\), and its positive-weight ideal is
\(\mathcal H^{\mot}_{>0}=\bigoplus_{w>0}\mathcal H^{\mot}_w\).
\end{definition}

\begin{definition}[Motivic coaction]
\label{def:vmpar-motivic-coaction}
\ClaimStatusDefinitional
The motivic coaction is Brown's graded coaction
\[
 \Delta^{\mot}\colon \mathcal H^{\mot}
 \longrightarrow \mathcal A^{\mathrm{dR}}\otimes\mathcal H^{\mot}.
\]
On \(q\in\mathbb Q\) it satisfies
\(\Delta^{\mot}(q)=1\otimes q\).  We call such a class Tate of
weight zero.
\end{definition}

\begin{definition}[Period map]
\label{def:vmpar-period-map}
\ClaimStatusDefinitional
The period homomorphism
\(\per\colon\mathcal H^{\mot}\to\mathbb R\)
sends a motivic iterated integral to its convergent numerical period
and restricts to the standard embedding on \(\mathbb Q\).
\end{definition}

\begin{definition}[Rational OPE package]
\label{def:vmpar-q-rational-ope}
\ClaimStatusDefinitional
Let \(F\) be a field of characteristic zero.  A strongly generated
vertex algebra has an \(F\)-rational OPE package when its singular OPE
coefficients, translation operator, vacuum functional, and invariant
pairing are defined over \(F\).  This package is local and belongs to
Beilinson level~\(1\).
\end{definition}

\begin{example}[Rational local families]
\label{ex:vmpar-q-rational-families}
The universal Virasoro vertex algebra has an
\(F=\mathbb Q(c)\)-rational OPE package:
\[
 T(z)T(w)\sim \frac{c/2}{(z-w)^4}
 +\frac{2T(w)}{(z-w)^2}+\frac{\partial T(w)}{z-w}.
\]
Universal affine vertex algebras and universal principal
\(\mathcal W\)-algebras carry analogous rational packages over their
parameter fields after the invariant form and strong generators have
been fixed.
\end{example}

\section{The all-arity motivic implication}
\label{sec:vmpar-main}

\begin{theorem}[Tate factorization of the Virasoro residue tower]
\label{thm:virasoro-s-r-motivic-purity-all-r}
\ClaimStatusConditional
\emph{Type signature:}
\textup{(Open quadrant, completed ordered residue-bar presentation,
Beilinson level~\(2\),
\(\mathsf H_{\mathrm{Vir}}+\mathsf H_{\mathrm{res}}
+\mathsf H_{\mathrm{mot}}+\mathsf H_{\mathrm{Tate}}\)).}
Let
\(S_r(\Vir_c;\mathsf H_{\mathrm{res}})\) be the scalar obtained from
the arity-\(r\) transferred Maurer--Cartan component.  If its motivic
lift factors through
\[
 \mathbb Q(c)(0)\hookrightarrow
 \mathcal H^{\mot}\otimes_{\mathbb Q}\mathbb Q(c),
\]
then, for every arity for which the residue package is defined,
\[
 S_r(\Vir_c;\mathsf H_{\mathrm{res}})\in\mathbb Q(c),
 \qquad
 \Delta^{\mot}S_r=1\otimes S_r.
\]
Thus the complete residue tower is motivically Tate of weight zero.
\end{theorem}

\begin{proof}
The hypothesis places the lifted scalar in the weight-zero summand
\(\mathbb Q(c)(0)\).  Definition~\ref{def:vmpar-motivic-coaction}
gives its coaction, and the period map returns the same rational
function.  The argument applies arity by arity.
\end{proof}

\begin{remark}[Logical relation with a formal recurrence]
\label{rem:vmpar-independence-from-riccati}
A rational Riccati recurrence proves rationality for the formal
sequence that it defines.  Theorem~\ref{thm:virasoro-s-r-motivic-purity-all-r}
concerns the geometric residue sequence.  The bridge between the two
is the comparison package
\(\mathsf H_{\mathrm{res}}+\mathsf H_{\mathrm{quad}}\) of
Chapter~\ref{ch:shadow-tower-higher-coefficients}.
\end{remark}

\section{Class-M algebras}
\label{sec:vmpar-class-m}

\begin{theorem}[Tate criterion for a rationally presented residue tower]
\label{thm:class-M-motivic-purity-algebras-with-Q-rational-OPE}
\ClaimStatusConditional
Let \(\cA\) carry an \(F\)-rational OPE package.  Assume
\(\mathsf H_{\mathrm{res}}(\cA)\), a motivic lift
\(\mathsf H_{\mathrm{mot}}(\cA)\), and a scalar contraction that
factors through \(F(0)\).  Then every defined scalar coefficient
\(S_r(\cA;\mathsf H_{\mathrm{res}})\) lies in \(F\) and has motivic
weight zero.
\end{theorem}

\begin{proof}
The contraction lands in \(F(0)\); the motivic coaction is therefore
the weight-zero coaction.  Period realization preserves the scalar.
\end{proof}

\begin{corollary}[Concrete rational families]
\label{cor:vmpar-concrete-Q-rational-families}
\ClaimStatusConditional
The criterion applies to a Virasoro, affine, or \(\mathcal W\)-family
once its ordered residue package, motivic realization, and Tate
factorization have been constructed over the displayed parameter
field.
\end{corollary}

\section{Where a positive motivic weight can enter}
\label{sec:vmpar-weight-of-mzv}

\begin{proposition}[Sources of positive motivic weight]
\label{prop:mzv-would-enter-at-what-weight}
\ClaimStatusProvedHere
For a residue construction with rational local OPE data, a
positive-weight motivic class can enter through a chain integral, a
regularized boundary value, an associator-dependent comparison, or a
period-valued contraction.  Its weight is the grading of that input in
\(\mathcal H^{\mot}\).  Rational algebraic operations preserve weight
zero.
\end{proposition}

\begin{proof}
Tensor products and rational linear combinations of weight-zero
classes remain in weight zero.  Consequently each positive-weight
summand is carried by one of the additional period-valued inputs of
the composite construction.
\end{proof}

\begin{remark}[Generator rigidity and residue geometry]
\label{rem:vmpar-virasoro-is-generator-rigid}
The Virasoro algebra has a single strong generator, so its local Ward
recursion is rigid.  The residue package still contains global chain
data.  This division identifies precisely where a motivic comparison
theorem acts.
\end{remark}

\section{Structural formulation}
\label{sec:vmpar-structural-reason}

\begin{theorem}[Structural Tate criterion]
\label{thm:structural-reason-for-purity}
\ClaimStatusConditional
Suppose a scalar invariant is represented by a composite
\[
 C_\bullet^{\mathrm{res}}(\cA)
 \xrightarrow{\ \mathcal T\ }
 H_\bullet^{\mathrm{res}}(\cA)
 \xrightarrow{\ \pi_{\mathrm{sc}}\ }F(0),
\]
where \(\mathcal T\) is the transferred Maurer--Cartan construction
and both arrows admit motivic lifts.  Then the scalar invariant is a
Tate period of weight zero.
\end{theorem}

\begin{proof}
The motivic composite factors through \(F(0)\), whose coaction is the
weight-zero coaction.
\end{proof}

\begin{remark}[Associator and residue integrals]
\label{rem:vmpar-associator-vs-bar}
An Etingof--Kazhdan comparison can contribute associator periods.  A
configuration-space residue model can contribute periods through its
own chains.  The package \(\mathsf H_{\mathrm{mot}}\) records both
sources, and \(\mathsf H_{\mathrm{Tate}}\) states their scalar
factorization.
\end{remark}

\begin{remark}[The GRT action on the Tate scalar line]
\label{rem:vmpar-grt-trivial-on-shadow}
When the motivic residue construction is \(\mathrm{GRT}_1\)-equivariant
and the final projection factors through \(F(0)\), the induced
\(\mathrm{GRT}_1\)-action on the scalar line is trivial.  The action on
the full chain model may retain its filtered associator directions.
\end{remark}

\section{Status of the all-arity claim}
\label{sec:vmpar-closure}

\begin{corollary}[All-arity Virasoro motivic status]
\label{cor:item-19-closed}
\ClaimStatusOpen
The all-arity Virasoro purity problem is reduced to the construction
of \(\mathsf H_{\mathrm{res}}\), \(\mathsf H_{\mathrm{mot}}\), and
\(\mathsf H_{\mathrm{Tate}}\), together with their compatibility with
the scalar projection.  The formal rational recurrence supplies an
independent candidate sequence after the package
\(\mathsf H_{\mathrm{quad}}\) has been added.
\end{corollary}

\section*{Epilogue}
\label{sec:vmpar-epilogue}

The local Virasoro algebra is rational.  The motivic theorem belongs
to the residue comparison.  This separation turns purity into a
finite chain-level construction: define the motivic residue model,
identify its scalar Tate line, and verify the transferred projection.
